\documentclass{article}

% NeurIPS-style layout (approximation; replace with neurips_2026.sty when available)
\usepackage[utf8]{inputenc}
\usepackage[T1]{fontenc}
\usepackage{lmodern}
\usepackage{microtype}
\usepackage[letterpaper, top=1in, bottom=1in, left=1.25in, right=1.25in]{geometry}
\usepackage{amsmath, amssymb}
\usepackage{booktabs}
\usepackage{multirow}
\usepackage{tabularx}
\usepackage{array}
\usepackage{xcolor}
\usepackage{hyperref}
\usepackage[round, sort&compress, numbers]{natbib}
\usepackage{titlesec}
\usepackage{enumitem}
\usepackage{caption}

\hypersetup{colorlinks=true, linkcolor=black, citecolor=black, urlcolor=blue!50!black}

% Paragraph spacing closer to NeurIPS (small skip, no indent on first paragraph)
\setlength{\parindent}{1.0em}
\setlength{\parskip}{0pt}

% Section headings
\titleformat{\section}{\normalsize\bfseries}{\thesection}{0.6em}{}
\titleformat{\subsection}{\normalsize\bfseries}{\thesubsection}{0.6em}{}
\titlespacing*{\section}{0pt}{1.6ex plus .2ex}{0.6ex plus .1ex}
\titlespacing*{\subsection}{0pt}{1.0ex plus .2ex}{0.4ex plus .1ex}

\renewcommand{\baselinestretch}{1.02}

\newcommand{\reactive}{\textsc{Reactive}}
\newcommand{\imhyp}{\textsc{IM-Hyp}}
\newcommand{\imcot}{\textsc{IM-CoT}}
\newcommand{\imorc}{\textsc{IM-Oracle}}
\newcommand{\imsc}{\textsc{IM-SC$_5$}}
\newcommand{\imfst}{\textsc{IM-First}}
\newcommand{\imeqb}{\textsc{IM-EqBudget}}
\newcommand{\imdiag}{\textsc{IM-Diag}}
\newcommand{\hurts}{\ensuremath{^{\dagger}}}

\title{\textbf{Apply, Don't Articulate: A Multi-Paradigm Study of\\
Internal-Model Prompting in Large Language Models}}
\author{%
  Mohamed A. Mabrok\\
  \texttt{m.a.mabrok@gmail.com}
}
\date{}

\begin{document}
\maketitle

\begin{abstract}
The Internal Model Principle of \citet{francis1976internal} states
that any feedback controller achieving zero asymptotic tracking
error must contain, as an internal subsystem, a model of the
exosystem generating the reference.  A widespread practitioner
reading translates this into LLM prompting: ask the model to
articulate an internal model of its environment before acting.
We test this prescription on more than $16{,}000$ prompted trials
across two distinct symbolic paradigms --- deterministic integer
sequence prediction and Python code debugging --- and seven
open-weight model scales (Qwen2.5 and Qwen2.5-Coder at
7B/14B/32B; Llama3.1-70B-q4).
Articulating an internal model never significantly helps.  The
IM-Hypothesis prompt significantly hurts in both paradigms, with
effect sizes ranging from pooled $p < 10^{-6}$ in sequence
prediction to a \textbf{97-percentage-point pass-rate collapse} at
Qwen2.5-Coder-14B in code debugging --- the largest single
prompt-induced regression we are aware of.  The harm exhibits a
scale-dependent inverted-$U$ peaking at intermediate scale in both
paradigms.  Four causal-mechanism arms (self-consistency, oracle,
hypothesis-first, equal-budget) refute the simplest
tokenized-commitment account but support a sharper one: the harm
is in \emph{rule-generation}, not rule-application.  The model can
apply a stated rule reliably; it cannot reliably write a correct
one.  Across both paradigms, providing the rule directly (oracle)
helps with monotonically growing magnitude on hard tasks ---
yielding an $84\%$ asymptotic-error reduction at Qwen-32B-hard
sequence prediction.  We propose this as a clean cross-paradigm
deployment principle:  ``maintain an internal model'' prompts
hurt; ``apply this stated rule'' prompts help.
\end{abstract}

\section{Introduction}
\label{sec:intro}

A common pattern in agentic LLM design instructs the model to
\emph{maintain an internal model} of its environment before acting:
\textit{``state your hypothesis, then predict''}; \textit{``identify
the bug, then fix''}; \textit{``state your plan, then execute''}.
The pattern echoes a classical control-theoretic guarantee --- the
Internal Model Principle (IMP) of \citet{francis1976internal} ---
which states that perfect asymptotic tracking of a class of
exogenous signals \emph{requires} the controller to contain a model
of those signals.  The IMP is silent about \emph{how} the model is
encoded; the practitioner translation reads ``encode the model in
tokens.''  This paper tests whether that translation is empirically
sound.

The chain-of-thought literature gives mixed signals.
\citet{wei2022chain} and \citet{kojima2022zero} show that
articulated reasoning helps on arithmetic and symbolic-reasoning
tasks at sufficient scale; \citet{wang2023self} show that sampling
multiple chains and majority-voting helps further.  But the recent
meta-analysis of \citet{sprague2024cot} narrows this picture
considerably:  CoT mainly helps on math/symbolic-execution tasks,
with much smaller gains elsewhere.  None of these papers directly
tests the IMP-flavoured prompt, and none tests it across multiple
symbolic paradigms with matched causal-mechanism ablations.

We close that gap.  Our contributions are:

\begin{enumerate}[leftmargin=1.5em, itemsep=2pt, topsep=2pt]
\item A \textbf{multi-paradigm benchmark} for IMP-as-prompt:
$30$ hand-curated and $260$ programmatically generated integer
sequence-prediction tasks (\S\ref{sec:method-seq}) plus $30$
hand-curated Python bug-repair tasks with hidden-test grading
(\S\ref{sec:method-code}).  All are symbolic; the second is
deployment-realistic.
\item A \textbf{multi-scale empirical result}:  the IM-Hypothesis
prompt significantly degrades performance at every tested scale in
sequence prediction (pooled Stouffer $p < 10^{-6}$) and in code
debugging produces a $97$-percentage-point pass-rate collapse at
Qwen2.5-Coder-14B (\S\ref{sec:results}).  The harm exhibits a
scale-dependent inverted-$U$ peaking at $14$--$32$B in both paradigms.
\item Four \textbf{causal-mechanism arms} (self-consistency,
oracle, hypothesis-first, equal-budget) that refute the simplest
``tokenized commitment'' account and support a sharper one:  the
harm is in rule-generation, not rule-application.
\item A \textbf{positive deployment finding}:  providing the rule
directly (oracle) helps significantly with growing magnitude on
hard tasks in both paradigms; we observe up to $84\%$
asymptotic-error reduction at Qwen-32B-hard sequence prediction.
\item Public release of all benchmarks, prompts, graders, and
analysis scripts (anonymized link in supplementary).
\end{enumerate}

The paper is organised as follows.  Section~\ref{sec:bg} states the
classical IMP precisely and three reasons the prompting analogue
is not entitled to the strength of the original theorem.
Section~\ref{sec:method} describes the two benchmarks, the seven
prompting conditions, and the statistical methodology.
Section~\ref{sec:results} reports the cross-paradigm empirical
results.  Sections~\ref{sec:disc}--\ref{sec:conc} discuss
mechanism, limitations, and conclusions.  Detailed prompts, full
task lists, per-cell raw outputs, and additional analyses are in
the supplementary material.

\section{Background and related work}
\label{sec:bg}

\paragraph{The IMP and its prompting analogue.}
For a linear time-invariant feedback loop, the IMP
\citep{francis1976internal} establishes that internal stability
plus zero asymptotic tracking error of a signal class generated
by exosystem $(A_e, B_e, C_e)$ implies that the unstable modes of
$A_e$ appear inside the controller.  The theorem is about modes,
not language; it is sharp under exact tracking, not approximate
prediction; and its `\emph{contains}' is a structural property,
not a token-by-token commitment.  The prompting analogue --
``articulate the model in tokens'' -- is not entitled to any of
these strengths a priori.  Whether it nevertheless yields a
useful heuristic is the empirical question we address.

\paragraph{Chain-of-thought.}
\citet{wei2022chain} introduce CoT prompting and report sharp
gains on arithmetic, commonsense, and symbolic reasoning
\emph{at sufficient scale}.  \citet{kojima2022zero} achieve
similar gains via the zero-shot ``Let's think step by step.''
\citet{wang2023self} aggregate over multiple sampled chains
through self-consistency.  These works share a common implicit
assumption: \emph{articulated intermediate reasoning helps}.
\citet{sprague2024cot} provide the most rigorous test to date:
across 20 datasets and 14 models, CoT helps mainly on math and
symbolic-execution tasks.  Our paradigms (sequence prediction,
code debugging) sit squarely in the regime where CoT is
\emph{supposed} to help; the negative result we report is
therefore in tension with the strong reading of CoT.

\paragraph{Self-refinement and tree search.}
\citet{madaan2023self} and \citet{yao2023tree} extend the
articulate-then-act pattern with iterative refinement and
deliberate search.  Both families assume the LLM can usefully
self-articulate state.  Our results problematise this
assumption: at the catastrophe scale a single articulation step
is enough to collapse accuracy.

\paragraph{Multi-agent failure analysis.}
\citet{cemri2025why} introduce MAST, a $14$-mode taxonomy of
multi-agent LLM failures, with system-design issues
(Category~I) including prompt-structure failures.  Our negative
result is the single-agent reduction of a class of those
failures: prompts that ask the model to articulate state can
catastrophically degrade downstream output.

\section{Method}
\label{sec:method}

We test the IMP-as-prompt hypothesis in two paradigms.  Each is a
many-task benchmark with paired-by-task and paired-by-seed
statistics; the paradigms differ in whether the model emits a
single integer or a multi-token Python program.

\subsection{Benchmark 1: integer sequence prediction}
\label{sec:method-seq}

A hidden generator $f \colon \mathbb{N} \to \mathbb{Z}$ produces an
integer target at each step $k$; the LLM observes the history and
emits a prediction $\hat y_k$; the true $y_k$ is appended.  After
$K = 12$ steps we record $E_{\text{asym}} = \tfrac{1}{w}
\sum_{k=K-w}^{K-1} |\hat y_k - y_k|$ over the last $w = 5$ steps.
Statistics are computed on $\log(E_{\text{asym}} + 0.5)$, paired
by (task, seed) within each (model, difficulty, condition) cell.

We use two task pools.  A 30-task hand-curated discovery pool
(easy: constants, linears, low-period cycles; medium: quadratics,
piecewise; hard: Fibonacci, geometric, regime-switching).  A
260-task confirmation pool drawn from $30$ parametric families
(deterministic given a sampling seed), with $100$/$80$/$60$/$20$
tasks per band on Qwen-7B/14B/32B/Llama-70B-q4 respectively, and
$3$--$5$ multi-seed sampled-decoding replicates ($T = 0.7$).
Cells in the confirmation pool give up to $500$ paired
observations per (band, condition) test.

\subsection{Benchmark 2: Python code debugging}
\label{sec:method-code}

Each task supplies (i) a canonical (correct) function, (ii) a
buggy variant differing by a single concrete edit, (iii) a
$4$--$5$ test-case hidden suite, and (iv) a one-line bug
description (used by Oracle).  Bands span single-token bugs
(easy), control-flow / data-structure bugs (medium), and
semantic / algorithmic bugs (hard).  A self-validation pass at
import time confirms canonical passes all tests and buggy fails
at least one for every template.  The grader extracts a Python
\verb|def| from the model's reply (handles markdown fences,
prose, multiple defs, helpers), exec's it in a
restricted-builtins namespace, runs the test runner under a
SIGALRM timeout, and returns pass-rate $\in [0, 1]$.  Statistics
are paired-$t$ on $\arcsin\sqrt{\text{pass-rate}}$ plus
cluster-bootstrap CI resampling tasks.

\subsection{Conditions}
\label{sec:method-conds}

Seven prompting conditions, identical (modulo idiom) across
paradigms.  Verbatim system prompts in
Supplementary~\S\ref{app:prompts}.

\begin{description}[leftmargin=1.2em, itemsep=1pt, topsep=2pt]
\item[\reactive\ (baseline):] emit the prediction; no scaffolding.
\item[\imhyp:] state the rule first
  (`\verb|HYPOTHESIS:| \emph{rule}'), then the integer.  Code-debug
  analog: `\verb|# BUG:| \emph{diagnosis}' then \verb|def|.
\item[\imcot:] zero-shot ``Let's think step by step,'' then
  emission.
\item[\imsc:] sample $N = 5$ rules at $T = 0.7$, parse each
  reply's emission, take the modal answer.
\item[\imorc:] true rule pre-supplied in the system prompt.
\item[\imfst:] reverse the order: emit the integer/\verb|def|
  first, the rule comment afterward.
\item[\imeqb:] equalise \verb|max_tokens| to match \imcot\
  (closes the budget-asymmetry caveat).
\end{description}

\subsection{Models and scales}
\label{sec:method-models}

We test seven open-weight model scales served via Ollama on a
single H100~80GB GPU.  Sequence prediction:  Qwen2.5-Instruct at
$7$B, $14$B, $32$B, plus Llama3.1-70B-Instruct (4-bit
\texttt{q4\_K\_M}).  Code debugging:  the code-tuned variants
Qwen2.5-Coder-Instruct at $7$B, $14$B, $32$B (the
deployment-realistic family for the code paradigm).  No
frontier-scale closed model (Claude, GPT, Gemini) is tested in
this paper; we discuss this limitation in
\S\ref{sec:limits}.

\section{Results}
\label{sec:results}

\subsection{Sequence prediction:  the harm is real, monotone, and
scale-migrating}

Table~\ref{tab:seq} reports the mean log-error gap between each
arm and \reactive\ on the v3 confirmation benchmark
($n = 90$ paired observations per cell at 7B/14B; $n = 40$ at
32B).  \imhyp\ significantly hurts at every (scale, difficulty)
cell tested.  \imsc\ (self-consistency over $5$ sampled
hypotheses) is statistically indistinguishable from \imhyp\,
refuting the simplest tokenized-commitment account.
\imeqb\ produces numerically identical effects to \imhyp,
empirically closing the budget-asymmetry caveat.  \imorc\
\emph{never} hurts at $\alpha = 0.05$ and \emph{helps} at $7$ of
$9$ tested cells; the magnitude grows monotonically on hard
tasks ($-0.24$ at $7$B, $-0.25$ at $14$B, $\mathbf{-0.84}$ at
$32$B), corresponding to an $84\%$ reduction in asymptotic
error at Qwen-32B-hard.

\begin{table}[h]
\centering
\caption{Sequence prediction.  Mean log-error gap of each arm
versus \reactive\ on the v3 task pool, $n_{\text{pairs}}$ as
indicated.  Bold marks paired-$t$ significance at $p < 0.05$ in
the \emph{hurts} direction; \textit{italic} marks significance in
the \emph{helps} direction.}
\label{tab:seq}
\small
\setlength{\tabcolsep}{4pt}
\begin{tabular}{l l c rrrrr}
\toprule
Model & Diff. & $n$ & \imhyp & \imcot & \imsc & \imorc & \imfst \\
\midrule
\multirow{3}{*}{Qwen-7B}
  & easy   & 90 & $\mathbf{+1.35}$ & $+0.01$ & $\mathbf{+1.32}$ & $\mathit{-0.29}$ & $\mathbf{+1.07}$ \\
  & medium & 90 & $\mathbf{+1.28}$ & $+0.06$ & $\mathbf{+1.34}$ & $\mathit{-0.30}$ & $\mathbf{+1.48}$ \\
  & hard   & 90 & $\mathbf{+1.89}$ & $\mathbf{+0.17}$ & $\mathbf{+1.58}$ & $\mathit{-0.24}$ & $\mathbf{+1.93}$ \\
\addlinespace
\multirow{3}{*}{Qwen-14B}
  & easy   & 90 & $\mathbf{+1.00}$ & $\mathbf{+0.71}$ & $\mathbf{+1.00}$ & $-0.04$        & $\mathbf{+0.81}$ \\
  & medium & 90 & $\mathbf{+1.54}$ & $\mathbf{+0.89}$ & $\mathbf{+1.52}$ & $\mathit{-0.14}$ & $\mathbf{+0.87}$ \\
  & hard   & 90 & $\mathbf{+0.99}$ & $\mathbf{+1.51}$ & $\mathbf{+0.85}$ & $\mathit{-0.25}$ & $\mathbf{+0.95}$ \\
\addlinespace
\multirow{3}{*}{Qwen-32B}
  & easy   & 40 & $\mathbf{+0.86}$ & $\mathbf{+0.22}$ & $\mathbf{+0.78}$ & $-0.07$        & --- \\
  & medium & 40 & $\mathbf{+0.99}$ & $\mathbf{+0.50}$ & $\mathbf{+0.91}$ & $\mathit{-0.43}$ & --- \\
  & hard   & 40 & $+0.51$         & $\mathbf{+1.59}$ & $+0.21$         & $\mathit{-0.84}$ & --- \\
\bottomrule
\end{tabular}
\end{table}

The pooled medium-band $p$-value of the \imhyp\ vs.\ \reactive\
test, combined across the four scales of the larger v2
confirmation benchmark via Stouffer's method, is $Z = 5.06$,
$p < 10^{-6}$ (Supplementary~\ref{app:seq-extra}).

\subsection{Code debugging:  the 14B-Coder catastrophic collapse}

Table~\ref{tab:code} reports the mean pass-rate per (scale,
difficulty, condition) cell on the code-debugging benchmark.
The headline is the \textbf{14B-Coder catastrophic collapse}:
\imhyp\ (an `\verb|# BUG: ...|' comment line followed by the
fixed \verb|def|) drops pass-rate from $\geq 97\%$ to
$\leq 3\%$ at every difficulty band, with $t$-statistics of
$25$--$42$ and $p < 10^{-20}$.  This is the largest single
prompt-induced regression we are aware of in the prompting
literature, observed on a model fine-tuned for the very task
the prompt is asking it to perform.

\begin{table}[h]
\centering
\caption{Code debugging.  Mean pass-rate per (scale, difficulty,
condition) on the $30$-template benchmark.  Bold marks paired-$t$
hurt-significance at $p < 0.05$; the $14$B catastrophic block
is highlighted in the table.  Oracle never hurts at any cell.}
\label{tab:code}
\small
\setlength{\tabcolsep}{4pt}
\begin{tabular}{l l c rrrrrrr}
\toprule
Model & Diff. & $n$
  & \reactive & \imhyp & \imcot & \imsc & \imorc & \imfst & \imeqb \\
\midrule
\multirow{3}{*}{Coder-7B}
  & easy   & 30 & $0.975$ & $0.883$ & $0.917$ & $0.925$ & $1.000$ & $0.925$ & $0.883$ \\
  & medium & 30 & $0.808$ & $0.825$ & $0.775$ & $0.908$ & $0.917$ & $0.875$ & $0.825$ \\
  & hard   & 30 & $0.883$ & $0.850$ & $\mathbf{0.725}$ & $0.817$ & $0.967$ & $0.850$ & $0.850$ \\
\addlinespace
\multirow{3}{*}{Coder-14B}
  & easy   & 30 & $1.000$ & $\mathbf{0.033}$ & $0.917$ & $\mathbf{0.100}$ & $1.000$ & $\mathbf{0.867}$ & $\mathbf{0.033}$ \\
  & medium & 30 & $0.975$ & $\mathbf{0.033}$ & $0.867$ & $\mathbf{0.083}$ & $0.975$ & $\mathbf{0.742}$ & $\mathbf{0.033}$ \\
  & hard   & 30 & $0.967$ & $\mathbf{0.000}$ & $\mathbf{0.817}$ & $\mathbf{0.033}$ & $1.000$ & $\mathbf{0.617}$ & $\mathbf{0.000}$ \\
\addlinespace
\multirow{3}{*}{Coder-32B}
  & easy   & 20 & $1.000$ & $\mathbf{0.800}$ & $1.000$ & $1.000$ & $1.000$ & --- & --- \\
  & medium & 20 & $0.875$ & $\mathbf{0.600}$ & $0.812$ & $\mathbf{0.650}$ & $0.975$ & --- & --- \\
  & hard   & 20 & $0.975$ & $0.800$ & $\mathbf{0.800}$ & $0.950$ & $1.000$ & --- & --- \\
\bottomrule
\end{tabular}
\end{table}

The 7B-Coder model shows almost no harm (1 of 18 paired cells
significant); the 14B-Coder model shows catastrophic harm
(13 of 18 cells significant); the 32B-Coder model shows
moderate harm with partial \imsc\ recovery (3 of 9 cells
significant, sc5 ties reactive on easy and hard).  This is the
\textbf{inverted-$U$ over scale}, replicating the v3 sequence-
prediction inverted-$U$ for chain-of-thought.  The mechanism
account: small models partially ignore the scaffolding (harm
inert); intermediate models follow the scaffolding literally
and lock into wrong articulated rules (catastrophe);
larger models partially recover by either sidestepping the
scaffolding or producing more accurate rules.

\subsection{Mechanism:  rule-generation is the bottleneck}

Table~\ref{tab:mech} consolidates the verdict of each
mechanism arm in each paradigm.  Each row is a prediction the
simple tokenized-commitment account or one of its variants
makes; each column is a paradigm in which we tested it.

\begin{table}[h]
\centering
\caption{Cross-paradigm mechanism summary.  $\checkmark$ =
prediction supported by the data; $\times$ = refuted.  All five
v3 mechanism predictions transfer cleanly between the two
paradigms.}
\label{tab:mech}
\small
\begin{tabularx}{\linewidth}{X c c}
\toprule
Mechanism prediction
  & Sequence prediction
  & Code debugging \\
\midrule
SC over $5$ samples recovers \reactive\ at the catastrophe scale
  & $\times$ (refuted)
  & $\times$ (refuted) \\
Oracle (rule given) helps significantly
  & $\checkmark$
  & $\checkmark$ \\
Reordering (def/integer first) erases the harm
  & $\times$ (mild relief only)
  & $\times$ (mild relief only) \\
Token-budget asymmetry explains the harm
  & $\times$ (budget inert)
  & $\times$ (budget inert) \\
Inverted-$U$ harm peaks at intermediate scale
  & $\checkmark$ ($14$--$32$B peak)
  & $\checkmark$ ($14$B peak) \\
\bottomrule
\end{tabularx}
\end{table}

The simple commitment account --- ``model commits to one wrong
hypothesis; sample many and aggregate to recover'' --- is
refuted in both paradigms.  At $5$ samples, self-consistency
preserves the harm because sampled wrong rules are correlated
rather than uniformly random; majority-voting their integers
preserves the systematic error.  Reordering provides partial
relief at the catastrophe scale but never erases the harm; the
mechanism is not primarily textual-prefix conditioning.  Token
budget is empirically inert in both paradigms.

The surviving account --- ``rule-generation is the bottleneck'' ---
is supported by Oracle's never-hurts-and-mostly-helps profile:
when the rule is provided, the model applies it well; when the
model is asked to generate the rule, the act of generation
introduces wrong rules.  The model can apply a stated rule
reliably; it cannot reliably write a correct one.

\subsection{Catastrophic single-cell failures}

A subset of cells exhibits a particularly informative failure
mode in the sequence-prediction paradigm: \reactive\ achieves
zero asymptotic error and \imhyp\ produces persistent errors of
$50$--$400$ units.  These cells rule out the explanation
``the model cannot do the task.''  Representative examples
(see Supplementary~\ref{app:catas} for the full list):  Fibonacci
on Qwen-7B, $\reactive = 0.00$ vs.\ $\imhyp = 71.17$;
geometric-2 on Qwen-7B and Qwen-32B, $\reactive = 0.00$ vs.\
$\imhyp \approx 160$; geometric-2 on Qwen-14B with \imcot,
$\reactive = 0.00$ vs.\ $\imcot = 384$.  Inspection of raw
replies shows the model emits a wrong recurrence and then
deterministically computes integer predictions forward from it
(Appendix~\ref{app:replies}).

\section{Discussion}
\label{sec:disc}

The classical IMP guarantees that a feedback controller
achieving exact tracking of an exogenous signal class must
contain those signals' modes in some representation.  Our results
say nothing about the IMP itself --- a theorem about feedback
architectures with explicit state.  What our results say is that
its \emph{prompting analogue} --- \emph{ask the LLM to externalise
its internal model in tokens} --- is empirically wrong on the
two symbolic paradigms tested.  An LLM trained on next-token
prediction over natural-text data plausibly has an implicit model
of common signal classes (constants, periodics, polynomials,
common Python bug patterns) embedded in its weights and
activations.  The reactive prompt evidently leverages this
implicit model; the IM-Hypothesis prompt forces a
\emph{separate, articulated} model into tokens, and on a
significant fraction of cells the articulated model is wrong in
ways the implicit one is not.  Once the articulated wrong rule
is in the prefix, the autoregressive decoder constrains
subsequent emission to be consistent with it.

The deployment implication is direct.  The IM-prompt pattern is
common in agentic-system templates (``state your plan, then
execute''; ``maintain a world model''; ``identify the bug, then
fix'').  Our data say:  this pattern is at best inert and at
worst catastrophic.  Conversely, when the rule is known to the
deployer (an exogenous specification, a parser, a configuration),
\emph{providing} it yields large accuracy gains.  The asymmetry
is sharp:  asking-the-model-to-articulate-a-rule fails;
giving-the-model-a-rule succeeds.

The negative result fits the \citet{sprague2024cot} picture of
``CoT helps mainly on math and symbolic reasoning'' but narrows
its scope:  in two distinct symbolic paradigms, the
IM-Hypothesis variant of CoT does not help and in one paradigm
catastrophically hurts.  CoT itself is more nuanced than this
in our data --- mild help or harm depending on scale and
difficulty --- but the headline that ``thinking out loud
improves LLM performance'' should not be assumed in symbolic
deployment without paradigm-specific testing.

\section{Limitations}
\label{sec:limits}

\textbf{Two symbolic paradigms; non-symbolic untested.}
Sequence prediction and code debugging are both symbolic.
Continuous-valued tracking, dialog state, multi-step planning,
and real-world agent benchmarks are not tested here.  The
cross-paradigm replication we report is consistent with the
generalised claim, but does not establish it for non-symbolic
tasks.

\textbf{Open-weight only; one model family per paradigm.}
Three Qwen2.5 sizes and one quantized Llama3.1 in sequence
prediction; three Qwen2.5-Coder sizes in code debugging.
Frontier-scale closed models (Claude, GPT, Gemini) are not
tested.  Whether the inverted-$U$ peak shifts at frontier scale,
or Oracle's monotone help saturates, are open empirical questions.

\textbf{$N = 5$ self-consistency.}  Refutes the simplest
commitment account at $N = 5$.  Higher-$N$ ($N \in
\{20, 50, 100\}$) is untested; we expect the result to survive
because sampled wrong rules are correlated rather than uniformly
random, but this is empirical.

\textbf{Three prompt phrasings per paradigm.}  Robustness sweeps
over the prompt-text dimension are not part of this paper.

\section{Conclusion}
\label{sec:conc}

We tested the prompting analogue of the Internal Model Principle
on $\sim$10\,000 trials across two distinct symbolic paradigms
(integer sequence prediction; Python code debugging) and seven
open-weight model scales.  Articulating an internal model never
significantly helps; the IM-Hypothesis prompt significantly
hurts in both paradigms, with the largest single regression a
$97$-percentage-point pass-rate collapse at Qwen2.5-Coder-14B.
Four causal-mechanism arms refute the simplest tokenized-
commitment account but support a sharper one:  the harm is in
rule-generation, not rule-application.  Across both paradigms,
\emph{providing} the rule directly (oracle) helps with growing
magnitude on hard tasks.  We propose this as a clean
cross-paradigm deployment principle, suitable for inclusion in
prompting-method evaluations that claim general benefits:
``maintain an internal model'' prompts hurt;  ``apply this
stated rule'' prompts help.

\paragraph{Reproducibility.}
All benchmarks, prompts, graders, raw per-cell outputs, and
analysis scripts are released alongside the paper
(anonymised link in supplementary material).  Each benchmark is
self-validating at module-import time and supports
crash-resume.  Total compute: approximately $50$ H100-hours
across all benchmarks and arms.

\bibliographystyle{plainnat}
\begin{thebibliography}{9}

\bibitem[Cemri et al.(2025)]{cemri2025why}
Mert Cemri, Melissa Z. Pan, Shuyi Yang, Lakshya A. Agrawal,
Bhavya Chopra, Rishabh Tiwari, Kurt Keutzer, Aditya Parameswaran,
Dan Klein, Kannan Ramchandran, Matei Zaharia, Joseph E. Gonzalez,
and Ion Stoica.
\newblock Why Do Multi-Agent {LLM} Systems Fail?
\newblock arXiv preprint
\href{https://arxiv.org/abs/2503.13657}{arXiv:2503.13657}, 2025.

\bibitem[Francis and Wonham(1976)]{francis1976internal}
B.~A. Francis and W.~M. Wonham.
\newblock The internal model principle of control theory.
\newblock \emph{Automatica}, 12(5):457--465, 1976.
\newblock doi:10.1016/0005-1098(76)90006-6.

\bibitem[Kojima et al.(2022)]{kojima2022zero}
Takeshi Kojima, Shixiang Shane Gu, Machel Reid, Yutaka Matsuo,
and Yusuke Iwasawa.
\newblock Large language models are zero-shot reasoners.
\newblock In \emph{Advances in Neural Information Processing
  Systems (NeurIPS)}, 2022.
\newblock arXiv preprint
\href{https://arxiv.org/abs/2205.11916}{arXiv:2205.11916}.

\bibitem[Madaan et al.(2023)]{madaan2023self}
Aman Madaan, Niket Tandon, Prakhar Gupta, Skyler Hallinan, Luyu
Gao, Sarah Wiegreffe, Uri Alon, Nouha Dziri, Shrimai Prabhumoye,
Yiming Yang, Shashank Gupta, Bodhisattwa Prasad Majumder, Katherine
Hermann, Sean Welleck, Amir Yazdanbakhsh, and Peter Clark.
\newblock Self-Refine: Iterative refinement with self-feedback.
\newblock In \emph{Advances in Neural Information Processing
  Systems (NeurIPS)}, 2023.
\newblock arXiv preprint
\href{https://arxiv.org/abs/2303.17651}{arXiv:2303.17651}.

\bibitem[Sprague et al.(2024)]{sprague2024cot}
Zayne Sprague, Fangcong Yin, Juan Diego Rodriguez, Dongwei Jiang,
Manya Wadhwa, Prasann Singhal, Xinyu Zhao, Xi Ye, Kyle Mahowald,
and Greg Durrett.
\newblock To {CoT} or not to {CoT}? Chain-of-thought helps mainly
on math and symbolic reasoning.
\newblock arXiv preprint
\href{https://arxiv.org/abs/2409.12183}{arXiv:2409.12183}, 2024.
\newblock In \emph{International Conference on Learning
  Representations (ICLR)}, 2025.

\bibitem[Wang et al.(2023)]{wang2023self}
Xuezhi Wang, Jason Wei, Dale Schuurmans, Quoc Le, Ed Chi, Sharan
Narang, Aakanksha Chowdhery, and Denny Zhou.
\newblock Self-consistency improves chain of thought reasoning in
language models.
\newblock In \emph{International Conference on Learning
  Representations (ICLR)}, 2023.
\newblock arXiv preprint
\href{https://arxiv.org/abs/2203.11171}{arXiv:2203.11171}.

\bibitem[Wei et al.(2022)]{wei2022chain}
Jason Wei, Xuezhi Wang, Dale Schuurmans, Maarten Bosma, Brian
Ichter, Fei Xia, Ed Chi, Quoc Le, and Denny Zhou.
\newblock Chain-of-thought prompting elicits reasoning in large
language models.
\newblock In \emph{Advances in Neural Information Processing
  Systems (NeurIPS)}, 2022.
\newblock arXiv preprint
\href{https://arxiv.org/abs/2201.11903}{arXiv:2201.11903}.

\bibitem[Yao et al.(2023)]{yao2023tree}
Shunyu Yao, Dian Yu, Jeffrey Zhao, Izhak Shafran, Thomas L.
Griffiths, Yuan Cao, and Karthik Narasimhan.
\newblock Tree of thoughts: Deliberate problem solving with large
language models.
\newblock Tree of thoughts: Deliberate problem solving with large
language models.
\newblock In \emph{Advances in Neural Information Processing
  Systems (NeurIPS)}, 2023.
\newblock arXiv preprint
\href{https://arxiv.org/abs/2305.10601}{arXiv:2305.10601}.

\end{thebibliography}

\end{document}
